\chapter{Computer Performance}

\section{Terminology}
\begin{itemize}[parsep=0.5em]
   \item $\text{Performance}_X = \frac{1}{\text{Execution time}_X}$
   \item Cycle time (seconds per cycle)
   \item CPU time / Execution time (seconds per program)
   \item CPU clock cycles (cycles per program)
   \item Instruction count (instructions per program)
   \item Clock rate (cycles per second)
   \item CPI (cycles per instruction)
   \item MIPS (millions of instructions per second)
\end{itemize}


\begin{itemize}
   \item \textbf{\textit{Response time}}
   : The time from when a request (or task) is submitted to when the system starts responding is the response time. Lower the response time indicates faster service.

   \item \textbf{\textit{Throughput}}
   : the total amount of work done in a given time.
\end{itemize}

\vspace{2em}
$\begin{aligned}
   & \text{Clock Cycles (No of Pulse required)} = \text{CPU Time} \times \text{Clock Rate} \\[0.5em]
   & \text{CPI} = \frac{\text{Clock Cycles}}{\text{Instruction Count}} = \frac{\text{CPU Time} \times \text{Clock Rate}}{\text{Instruction Count}} \\[0.5em]
   & \text{Execution time} = \text{Clock cycles} \times \text{Clock cycle time} \\[0.5em]
   & \text{Execution time} = \text{Instruction count} \times \text{CPI} \times \text{Clock cycle time}
\end{aligned}$




\pagebreak
\section{Numerical Problems}
\vspace{1em}
\textbf{\textit{\#Question-1}:}\\
A designer building a new machine B, which must run a specific program in 6 seconds. The current machine, A, completes the same program in 10 seconds, using a 2.5 GHz clock rate. The designer plans to increase the clock rate to improve performance. However, this design change will cause Machine B to require 1.2 times as many clock cycles as Machine A for the same program. Based on this information, calculate the clock rate that Machine B must have.

\textbf{\textit{Solution}:}
\begin{multicols}{2}
   \noindent
   \raggedcolumns
   Now,\\
   For Machine A,\\
   $\begin{aligned}
      \text{Clock Cycles } &= \text{CPU Time } \times \text{Clock Rate} \\
      &= \text{10} \times \text{2.5} \times \text{10}^{\text{9}} 
      = \text{2.5} \times \text{10}^{\text{10}} \text{ cycles}
   \end{aligned}$

   For Machine B,\\
   $\begin{aligned}
      \text{Clock Cycles } &= 1.2 \times \text{Machine A clock cycles} \\
      &= \text{1.2} \times \text{2.5} \times \text{10}^{\text{10}} 
      = \text{3} \times \text{10}^{\text{10}} \text{ cycles}
   \end{aligned}$

   $\begin{aligned}
      \therefore \text{CPU Time} &= \frac{\text{Clock Cycles}}{\text{Clock Rate}} \\[0.5em]
     \Rightarrow \text{Clock Rate} &= \frac{\text{Clock Cycles}}{\text{CPU Time}} \\
       &= \frac{\text{3} \times \text{10}^{\text{10}}}{\text{6}} \\
       &= \text{5} \times \text{10}^{\text{9}} \text{Hz}
       = \text{5} \text{GHz} \\
   \end{aligned}$
   \columnbreak
   Given,

   For Machine A,\\
   Execution time = 10 seconds\\
   Clock rate = 2.5 GHz = 2.5 $\times$ $\text{10}^{\text{9}}$ Hz

   For Machine B,\\
   Execution time = 6 seconds\\
   Clock cycles = 1.2 $\times$ Machine A clock cycles
\end{multicols}



\textbf{\textit{\#Question-2}:}\\
Suppose Machine $X$ needs 5 seconds to run a program and $Y$ needs 3 seconds to run the same program and Clock rate is 3 GHz. Machine $Y$ to require 2.5 times as many clock cycles as Machine X for the same program. What will be the clock rate?

\textbf{\textit{Solution}:}
\begin{multicols}{2}
   \noindent
   \raggedcolumns
   Now,

   For Machine X,\\[0.5em]
   $\begin{aligned}
      \text{Clock cycles} &= \text{CPU time} \times \text{Clock rate}\\
      &= \text{5} \times \text{3} \times \text{10}^{\text{9}}
      = \text{15} \times \text{10}^{\text{9}} \text{ cycles}
   \end{aligned}$

   For Machine Y,\\[0.5em]
   $\begin{aligned}
      \text{Clock cycles} &= \text{2.5} \times \text{Machine A clock cycles}\\
      &= \text{2.5} \times \text{15} \times \text{10}^{\text{9}}
      = \text{3.75} \times \text{10}^{\text{10}} \text{ cycles}
   \end{aligned}$

   $\begin{aligned}
     \therefore \text{CPU Time} &= \frac{\text{Clock Cycles}}{\text{Clock Rate}} \\[0.5em]
     \Rightarrow \text{Clock Rate} &= \frac{\text{Clock Cycles}}{\text{CPU Time}} \\
       &= \frac{\text{3.75} \times \text{10}^{\text{10}}}{\text{3}} \\
       &= \text{1.25} \times \text{10}^{\text{10}} \text{Hz}
       = \text{12.5} \text{GHz} \\
   \end{aligned}$


   \columnbreak
   Given,

   Machine X,\\
   Execution time = 5 seconds\\
   Clock rate = 3 GHz = 3 $\times$ $\text{10}^{\text{9}}$ Hz

   Machine Y,\\
   Execution time = 3 seconds\\
\end{multicols}


\pagebreak
\textbf{\textit{\#Question-3}:}\\
A compiler designer is trying to decide between two code sequences for a particular machine. Based on the hardware implementation, there are three different classes of instructions: Class A, Class B, and Class C, and they require 1, 2 and 3 cycles (respectively). The first code sequence has 5 instructions: 2 of A, 1 of B, and 2 of C The second sequence has 6 instructions: 4 of A, 1 of B, and 1 of C. Which sequence will be faster? How much? What is the CPI for each sequence?

\textbf{\textit{Solution}:}
\begin{multicols}{2}
   \noindent
   \raggedcolumns
   Now,\\[0.5em]
   For sequence-1,\\[0.5em]
   $\begin{aligned}
      \text{CPI} &= \frac{\text{Totol CPU Cycles}}{\text{Instruction Count}} \\[0.5em]
      &= \frac{(2 \times 1) + (1 \times 2) + (2 \times 3)}{5}\\
      &= \frac{10}{5} = 2 \text{ cycles per instruction}
   \end{aligned}$

   For sequence-2,\\[0.5em]
   $\begin{aligned}
      \text{CPI} &= \frac{\text{Totol CPU Cycles}}{\text{Instruction Count}} \\[0.5em]
      &= \frac{(4 \times 1) + (1 \times 2) + (1 \times 3)}{6}\\
      &= \frac{9}{6} = 1.5 \text{ cycles per instruction}
   \end{aligned}$
   

   Now,\\
   Let, Clock cycle time $=$ 1$ns$ \\[0.5em]
   $\begin{aligned}
      \therefore \text{ET}_{1} &= \text{Total cycles} \times \text{Clock cycle time}
      = \text{10} \times \text{1}
      \quad \ldots \quad (1) \\
      % 
      \therefore \text{ET}_{2} &= \text{Total cycles} \times \text{Clock cycle time}
      = \text{9} \times \text{1}
      \quad \ldots \quad (2)
   \end{aligned}$
   
   $(2) \div (1) \ \Rightarrow
   \begin{aligned}
   \frac{\text{ET}_{2}}{\text{ET}_{1}}
      = \frac{9}{10}  = 0.9 \\
   \end{aligned}$
   
   \vspace{0.5em}
   $\therefore \text{Sequence 2 is 0.9 times faster than sequence 1.}$

   \columnbreak
   Given,\\[0.5em]
   Cycles required for,\\
   Class A = 1\\
   Class B = 2\\
   Class C = 3

   For sequence-1,\\
   Instructions: 2 of A, 1 of B, and 2 of C = 5 cycles

   For sequence-2,\\
   Instructions: 4 of A, 1 of B, and 1 of C = 6 cycles
\end{multicols}
